\chapter{Μεθοδολογία Αξιολόγησης}

\section{Σχεδιασμός Αξιολόγησης Μαθησιακού Περιβάλλοντος (\en{IEQ} \& Ηχητική Άνεση)}
Βασικός άξονας αξιολόγησης αποτέλεσε η μέτρηση της ακουστικής άνεσης (\en{Noise Level}) στις σχολικές αίθουσες, με
στόχο τη διατήρηση του θορύβου υποβάθρου σε επίπεδα ευθυγραμμισμένα με τις οδηγίες του Παγκόσμιου Οργανισμού
Υγείας (\en{World Health Organization, WHO}), δηλαδή $\leq$ \texttt{35 dBA} για κενή αίθουσα
\cite{WHO2015SchoolEnvironment, Mercugliano2025, MogasRecalde2021}. Η μεθοδολογία περιέλαβε την καταγραφή του δείκτη \texttt{LAeq,10s} μέσω των
μικροφώνων \en{INMP441}. Ως μετρικές επιτυχίας ορίστηκαν: (α) ο αριθμός υπερβάσεων του ορίου ανά διδακτική ώρα και
(β) η συνολική διάρκεια των υπερβάσεων αυτών.

Η καταγραφή πραγματοποιήθηκε σε δύο διακριτές φάσεις:
\begin{itemize}
  \item \textbf{Περίοδος Αναφοράς (\en{baseline}):} Χωρίς οπτική ανατροφοδότηση.
  \item \textbf{Περίοδος Βελτιστοποίησης:} Με αυτόματη ενεργοποίηση προειδοποιητικής λυχνίας όταν παρατηρείται
  υπέρβαση του ορίου.
\end{itemize}

Η συλλογή και επεξεργασία των δεδομένων ακολούθησε ένα αυστηρό πρωτόκολλο, το οποίο περιέλαβε κριτήρια ένταξης,
διαχείριση ακραίων τιμών (\en{outliers}) και πολιτική συμπλήρωσης κενών δεδομένων (\en{missing data policy}). Στον
Πίνακα~\ref{tab:methodology_protocol} συνοψίζονται οι βασικές παράμετροι του πειραματικού σχεδιασμού.

\begin{table}[htbp]
\centering
\caption{Πρωτόκολλο Μεθοδολογίας και Επεξεργασίας Δεδομένων}
\label{tab:methodology_protocol}
\begin{tabular}{@{}p{4cm}p{10cm}@{}}
\toprule
\textbf{Παράμετρος} & \textbf{Περιγραφή / Κανόνας} \\
\midrule
\textbf{Περίοδοι Σύγκρισης} & Περίοδος Βάσης (\en{Baseline}): Σεπ 2024 -- Φεβ 2025. Περίοδος Βελτιστοποίησης
(\en{DT}): Μαρ 2025 -- Νοε 2025. \\
\textbf{Διαχείριση Outliers} & Αφαίρεση τιμών θορύβου \(> 100\) \texttt{dBA} (σφάλμα αισθητήρα) και μηδενικών
καταγραφών ισχύος (\texttt{0 W}) κατά τη διάρκεια σχολικού ωραρίου (\en{IQR filter}). \\
\textbf{Missing Data Policy} & Γραμμική παρεμβολή (\en{linear interpolation}) για κενά δεδομένων μικρότερα της
\texttt{1} ώρας. Εξαίρεση ολόκληρης της ημέρας για κενά \(> 1\) ώρας. \\
\textbf{Αριθμός Αιθουσών} & Τρεις (\texttt{3}) αίθουσες διδασκαλίας και ένα (\texttt{1}) κυλικείο. Αξιολόγηση μόνο
κατά τις ώρες μαθημάτων (\texttt{08:00--14:00}). \\
\bottomrule
\end{tabular}
\end{table}

\section{Μεθοδολογία Ενεργειακής Σύγκρισης}
Η αξιολόγηση βασίστηκε σε πραγματικά δεδομένα από λογαριασμούς ενέργειας. Για τη συγκριτική ποσοτικοποίηση
χρησιμοποιήθηκαν οι διαθέσιμες περίοδοι που παρατίθενται στο Παράρτημα Β: Περίοδος Βάσης
(Σεπτέμβριος 2024 -- Φεβρουάριος 2025) και Περίοδος Βελτιστοποίησης (Μάρτιος 2025 -- Νοέμβριος 2025). Ως βασική
μετρική επιλέχθηκε η Μέση Ημερήσια Κατανάλωση Ενέργειας (\texttt{kWh/day}). Παράλληλα, σχεδιάστηκε σύγκριση
\en{Year-over-Year (YoY)} για ταυτόσημους μήνες.

\section{Σχεδιασμός Τεχνοοικονομικής Μελέτης (Μοντέλο \en{ROI})}
Αναπτύχθηκε μοντέλο Υπολογισμού Επιστροφής Επένδυσης (\en{Return on Investment, ROI Calculator}), το οποίο
τροφοδοτήθηκε από το κόστος εξοπλισμού (\en{Bill of Materials}) και τα οικονομικά οφέλη που προκύπτουν από τη
μείωση του κόστους ηλεκτρικής ενέργειας και του λειτουργικού κόστους.
